\section{Gestion de l'énergie}
\label{sec:soft_power}

La régulation, la charge et la protection contre la décharge profonde sont gérées par des circuits autonomes décrits au chapitre précédent, échappant à tout contrôle direct du firmware. L'intervention logicielle se limite à cinq actions : le maintien sous tension, la surveillance énergétique, l'arbitrage du connecteur USB, la réduction de consommation et l'extinction contrôlée. Ces responsabilités sont regroupées au sein d'un service d'énergie dépourvu de tâche permanente. Son exécution intervient une fois sur vingt dans le contexte de la tâche de maintenance, ce qui établit sa période à deux secondes.

\subsection{Maintien sous tension et observation}

Au démarrage, le rail d'alimentation n'est maintenu que par le bouton de mise en marche ou par la présence de l'USB. La toute première instruction exécutée par le firmware met donc le signal EnableReg à l'état haut, faute de quoi le relâchement du bouton éteindrait la carte avant même la fin de l'initialisation.

À chaque cycle, le service interroge la jauge de charge sur le bus I\textsuperscript{2}C et relit le signal INOKB du chargeur via l'expanseur, signal à drain ouvert par lequel le chargeur avertit la présence d'une source valide à son entrée~\cite{analog_MAX77757}. Il en tient à jour un relevé unique qui réunit le \gls{soc}, la tension, le courant, le temps de charge et décharge restant estimé et la présence d'une alimentation externe. Ce relevé ne nécessite aucun verrou. Une lecture concurrente observera dans le pire des cas une ancienne valeur de cycle, ce qui reste acceptable pour l'affichage.

Deux mécanismes de repli protègent l'intégrité du relevé. Un échec de lecture de la jauge invalide les données et interrompt le cycle sans prise de décision, évitant ainsi toute action basée sur une mesure douteuse. Un échec de lecture de l'expanseur conserve en revanche l'état précédent du signal INOKB pour prévenir un débranchement fictif qui relancerait inutilement la séquence d'arbitrage et priverait la console UART de son port.

Ce relevé est enfin traduit en un indicateur de charge à trois niveaux : normal, bas et critique. Les seuils de 15\,\% et 5\,\% sont des réglages fixes appliqués à l'état de charge corrigé fourni par la jauge~\cite{analog_MAX17260}. La détection de charge en cours repose sur un courant entrant supérieur à 10~mA, seuil choisi pour distinguer une charge réelle du bruit de mesure au repos.

\subsection{Arbitrage du connecteur USB}

Le connecteur USB-C assure de manière exclusive soit la charge soit la console série car ses lignes de données ne peuvent être reliées qu'à un seul circuit à la fois. Cet arbitrage n'affecte pas la charge depuis une source Type-C moderne qui lit le courant disponible sur les broches d'alimentation sans utiliser les lignes de données~\cite{analog_MAX77757}. En revanche une source legacy reliée par un câble USB-A vers USB-C ne peut annoncer son courant disponible que via la négociation \gls{bc12} qui nécessite l'accès aux lignes de données~\cite{usb_bc12}. Sans cet arbitrage la limite d'entrée du chargeur resterait figée à sa valeur par défaut même face à une source capable de fournir davantage.

La séquence retenue suit la recommandation du constructeur en confiant d'abord les lignes de données au chargeur puis en attendant que le signal INOKB indique une source valide avant de les rendre au pont série~\cite{analog_MAX77757}. Cette opération dure plus d'une seconde et est déléguée à une tâche éphémère. Deux durées encadrent cette attente dont la première couvre le pire cas de détection estimé à 994~ms pour les trois étapes successives~\cite{analog_MAX77757}. La valeur retenue de 1200~ms inclut une marge de sécurité car rendre les lignes plus tôt interromprait la négociation en cours. La seconde durée agit comme garde-fou en abandonnant la tâche après cinq secondes si INOKB ne s'active jamais.

La détection n'ayant lieu qu'au branchement le cycle périodique mémorise l'état précédent de l'alimentation externe pour relancer l'arbitrage sur un front montant et ainsi identifier toute source connectée après le démarrage. Cette période de deux secondes implique deux limites assumées dont un délai potentiel avant la reconnaissance du branchement et l'impossibilité de détecter un débranchement suivi d'un rebranchement plus rapide que cette période. Deux mécanismes de protection complètent le système en ignorant toute seconde demande d'arbitrage tant que la première est en cours et en ramenant immédiatement les lignes sur le pont série en cas d'échec de création de la tâche. Cette règle garantit qu'aucune défaillance ne bloque les lignes USB sur le chargeur et conserve une console série utilisable même si la jauge cesse de répondre au prix d'un courant de charge plafonné.

\subsection{Ce que le firmware peut économiser}

Les modes de sommeil du microcontrôleur sont incompatibles avec la fonction de l'appareil. Le sommeil léger suspend les horloges du bus \gls{i2s} et interromprait la lecture audio filaire. Il exigerait en plus un quartz externe de 32,768~kHz pour maintenir le Bluetooth actif, quartz absent de la carte, ce qui le rend inopérant dans ce cas~\cite{ESP32_IDF_pm}. Le sommeil profond, qui perd le contexte d'exécution, est lui aussi exclu puisque l'appareil doit continuer à jouer de façon continue. L'économie possible ne porte donc que sur la façon dont le processeur travaille éveillé, et sur la décision de l'éteindre complètement.

La gestion dynamique de fréquence, ou \gls{dfs}, a été testée puis abandonnée. Configurée entre 80 et 160~MHz, elle provoquait des blocages aléatoires du système. L'analyse des \gls{coredump} a désigné le \gls{iwdt}, décrit à la section~\ref{sec:soft_robustness}, qui se déclenchait sur le second cœur pendant la mise en attente du module de gestion d'énergie~\cite{ESP32_IDF_wdt}. Cette signature est documentée par d'autres rapports publics~\cite{espressif_issue_2542, esp32_forum_wdt_pm}, et la documentation officielle reconnaît que la gestion d'énergie augmente la latence d'interruption~\cite{ESP32_IDF_pm}. Le détail figure en annexe~\ref{ann:coredump}. La fréquence a donc été figée à 160~MHz, sans autre blocage depuis. La cause précise n'a pas été creusée davantage, le code concerné appartenant au framework propriétaire d'Espressif. La stabilité prime ici sur l'économie d'énergie.

Trois leviers restent effectifs. Le modèle événementiel de la section~\ref{sec:archi} ne consomme aucun cycle au repos et constitue la principale économie. Les deux autres sont des politiques temporisées, l'extinction automatique de l'écran (section~\ref{sec:ui_verrouillage}) et de l'appareil, réglables par l'utilisateur et sauvegardées en \gls{nvs}, avec des valeurs par défaut de 30~secondes et 5~minutes. Une valeur nulle désactive ces options. Ces réglages sont relus une seule fois puis gardés en mémoire, pour éviter un accès à la flash à chaque cycle.

\subsection{Extinction ordonnée et arrêt automatique}

Trois voies conduisent à l'extinction, la commande explicite de l'utilisateur, l'atteinte du seuil critique de batterie et l'expiration du délai d'inactivité. La figure~\ref{extinction.drawio} détaille l'enchaînement et les gardes de chaque voie.

\fig[H, width=1\textwidth]{Séquence d'extinction. Diagramme créé avec l'aide de Claude Opus 4.8.}{extinction.drawio}

Les deux voies automatiques sont désactivées en présence d'une alimentation externe. Il ne s'agit pas d'un choix mais d'un constat, l'USB branché maintient le régulateur activé indépendamment du signal EnableReg, si bien qu'une extinction ne ferait qu'arrêter le firmware sur un appareil resté alimenté. L'extinction sur batterie critique est de plus suspendue pendant une campagne de mesure d'autonomie, ce service devant lui-même détecter le seuil et écrire son dernier relevé avant de demander l'extinction.

La condition d'inactivité ne tient compte que des appuis de boutons, la rotation du potentiomètre de volume ne la réinitialisant pas. Cette limite n'a de portée qu'en pause, la lecture interdisant de toute façon l'arrêt.

Le service d'énergie n'a aucune connaissance de la chaîne audio, l'ordre d'arrêt est déposé dans la file de la tâche audio plutôt que transmis directement. Cette inversion de dépendance préserve la règle de la sectiozn~\ref{sec:archi} et laisse à la tâche audio le temps de s'arrêter correctement. La position de lecture est enregistrée sous forme de chemin et non de rang dans la liste, ce qui la rend insensible à une réindexation de la carte. Un marqueur d'extinction volontaire est écrit en \gls{nvs} en tout dernier, étape décrite plus précisément à la section \ref{sec:controlled_extinction}, une fois qu'aucune étape ne peut plus échouer.

L'extinction sur batterie critique emprunte un chemin plus court, le service relâchant le rail sans arrêter la lecture ni démonter le système de fichiers, la carte n'étant utilisée qu'en lecture. Le marqueur est en revanche écrit, sans quoi la coupure serait relue comme une extinction subie.

L’allumage du lecteur s’effectue exclusivement par voie matérielle, il suffit de maintenir le bouton A enfoncé pendant une à deux secondes. Cette durée n’est pas arbitraire, mais imposée par la séquence d’initialisation de l’ESP32.